\section{Cosmic and Positive Law}

With the elections in the USA in full swing, all the talking heads are claiming to be ``true conservatives''. The caterwauling means nothing, since they seem to believe, and with passionate intensity, what no previous conservative had ever believed. Unable to think in depth, people can only think in terms of so called ``issues'', aligning them on either liberal or conservative side. Unfortunately, the process is entirely arbitrary, based on nothing but the bigoted opinions of the speaker.

However, if God is the Moral Governor of the Universe, and the Logos, through whom all things were created, is the cosmic law, then we indeed have a sound basis to evaluate the ``issues''. Positive law must conform to the divine order in the vertical dimension. Then the horizontal categories of left-right becomes relativized.

\paragraph{Origins of Law}


In the traditional, vertical view --- that encompasses various strata of being --- we note that there are laws applicable at each layer. At the material layers, the laws are quantitative and precise, allowing no deviation. At the biological level, the actions of an animal are freer, but still subject to the laws of fear, hunger, and sex. When we reach the moral law, there is complete freedom. This freedom, or lack of compulsion, has led many to the false conclusion that there is no moral law or that it is merely conventional.

The truth of the matter is that we are to follow the cosmic law freely, out of love. No one wants love from someone who is coerced to give it. Tradition shows us three levels of law:

\begin{itemize}


\item Cosmic or Divine Law

\item Natural Law

\item Positive Law
\end{itemize}


Writing from a Thomist perspective, \textbf{Valentin Tomberg}, in Foundations of International Law as Humanity's Law, draws out the logical consequences of these ideas. The driving force or content behind these laws are respectively Love, Justice, and Power.

\textbf{Divine Law}: This is the law of the Spiritual Authority. It expresses the ideal of law, which is \textbf{Love}, as the intuition of the moral world order.

\textbf{Natural Law}: Everything that is created has its own nature.

\textbf{Justice} demands that everything be given its due, hence it is the basis for natural law. Only by understanding man's true nature, is it possible to understand the requirements of justice.

\textbf{Positive Law}: This is man-made, or legislative, laws and, as such, is the domain of the temporal \textbf{Power}. Positive law has three conditions:

\begin{itemize}


\item It must be necessary or useful for the general welfare

\item It must not contradict any obligations of a superior nature

\item It must physically and morally be possible
\end{itemize}


The first condition expresses the purpose of positive law. The second condition asserts that the positive law must be oriented toward justice and love. The third condition requires the virtue of Prudence. That is, the law must be Reasonable, and what is reasonable at one time for one nation may be different in other circumstances. The cost of enforcing a law may exceed any moral benefits arising from it. The Ethical State will provide the atmosphere for individual freedom and personal fulfillment. But the degree of that depends on the emotional and spiritual maturity of a people.

With these principles in mind, we can see how to approach economic issues.

\begin{enumerate}


\item The purpose of the economic system is to deliver the goods necessary for life: food, shelter, education, healthcare, leisure, etc. This includes workers as well as those incapable of working.

\item The money power must be subservient to the just laws of the temporal and spiritual authorities. Undue influence of the money class over the political class must be resisted.

\item Workers must be paid a reasonable or living wage. This area requires sound judgment to arbitrate competing demands.
\end{enumerate}


Thomas Aquinas' thoughts on these matters are summarized in \textit{The Governance of Princes}\footnote{\url{http://www.gornahoor.net/?p=69}}.

\paragraph{Law of Love}


Besides money and the economy, laws pertaining to sex in various manifestations cause a lot of contention. \textbf{Boris Mouravieff} draws on the Eastern Orthodox tradition for a deeper understanding of sexual energy or \emph{eros}. This energy manifests on three levels:

\begin{itemize}


\item \textbf{Spirit}: Love as creative energy


\item \textbf{Soul}: emotional and psychical support for Love


\item \textbf{Corporeal}: carnal love

\end{itemize}


Through the sublimation of this energy, it passes from procreation to creation. In brief, Mouravieff brings up several points: Woman, the passive force in conception becomes the active force in the act of creation. The passage from animal love to truly human Love is on the way towards the objective Love in which the Sermon on the Mount invites us to participate. This objective Love gives us a new enlarged and refined vision of more subtle qualities of Love which come close to the first impulse of Creation.

Not everyone, however --- not even most people --- are aware of the higher purposes of eros. Even most who are religious, and hence knowledgeable about the teachings of Love, don't really know how to transcend and sublimate sexual desire. Hence, there is disorder in social relations. There may have been a time when the teachings of the spiritual authority were accepted by faith, but that is hardly possible now. Since we are interested here in understanding the positive law on issues related to sex, we can bypass for now its higher manifestations and focus just on carnal love. Carnal love develops in four stages:

\begin{enumerate}


\item A general atmosphere of sexual attraction

\item The particular sexual attraction of a couple

\item Sexual union

\item Orgasm
\end{enumerate}


Deformations can occur at any stage. First of all, sexual attraction in itself is promiscuous and undifferentiated, i.e., not oriented toward anything in particular: for example, a man's desire may be for his wife's vagina, a prostitute's mouth, Rosie Palm and her five sisters, or even his buddy's rectum. This is similar perhaps to the concept of \emph{polymorphous perversity} in psychoanalysis. It is not an accident that revolutionary movements see this as the goal rather than the beginning of sexual maturity.

The next phase in the normative development of the sexual function is the attraction between a man and a woman. Once again, the strength of the attraction can lead to less than ideal results. Ultimately, there is one right man for each woman, although there may be many who are of a compatible type. This is why marriage in indissoluble. Unfortunately, the possibility of mismatches is quite high and this leads to a lot of unhappiness in life.

Sexual union and orgasm then result in reproduction, which is the biological justification for sexual activity. However, there is an excess of sexual energy beyond what is necessary for the continuance of the human race. If this energy is mastered, it can be applied to one's spiritual development. In most cases it is just frittered away in carnal love unrelated to reproductive ends. Through the power of the imagination, the mass of people seek even to artificially stimulate the force of sexual attraction. The entertainment industry, as well as commercial advertising, capitalizes on this need for their own benefit. This keeps the populace in bondage.

Attraction and union explain why there is a differentiation of the sexes not just on the corporeal level, but also on the psychical. However, in the vertical direction, through the sublimation of the sexual impulse, there can result the alchemical marriage of opposites. We can't discuss this at this point since our goal is to describe the foundations of law. Note how the temporal power, when it is divorced from any higher spiritual support, will encourage sexual promiscuity and dedifferentiation in order to control the populace.

Once the role of sex in the cosmic order is understood, the issues can be understood on a rational basis. Laws that seem to be arbitrary or stultifying will appear more reasonable. First of all, the law needs to encourage the linkage between reproduction and sexual activity. Otherwise, there is a demonstrable and measurable effect on the general welfare. Populations that fail to maintain themselves quantitatively are doomed to extinction. The turmoil in Europe, for example, around migration issues, would never have occurred if Europeans had been managing a healthy population growth. The rise of so-called ``right wing'' groups in Europe will not amount to much, unless they encourage their adherents to have large families. There is a real conflict here behind the desires of people to have a small family versus the general welfare.

The second function of the sexual impulse is to lead to higher manifestations of Love, which is the goal of the Ethical state. Mouravieff points out that there are two kinds of imagination:

\begin{enumerate}


\item \textbf{Creative imagination}: This is awake and constructive. It is this divine force which distinguishes men from beasts: it is an active force.

\item \textbf{Dreamlike imagination}: This is somnolent and is also found to a certain degree in animals: it is a passive force.
\end{enumerate}


Pornography, therefore, is undesirable, not out of prudery, but because it fosters the dreamlike imagination, or daydreaming. Although Mouravieff does not mention it, the drug experience, e.g., from opiates and psychedelics, is similar. The user gets totally absorbed in a waking dream sequence. It is curious that ``sophisticated'' people today are so dedicated to the sex and drug experience. They are not immune to its attractions. Mouravieff writes:

\begin{quotex}
The dreamlike imagination, the `dream of the sleeping serpent', produces a hypnotic effect on man, keeping him in the state in which the vast majority of humans pass their lives.
\end{quotex}


That is why sophisticates are so cocksure of their socio-political views: there can be little doubt about the contents of one's own imagination, no matter how degenerate it may be. However, the creative imagination is riskier and less certain. Nevertheless, the spiritual authority of a community needs to discourage the one and encourage the other, through the positive law if needs be.

\paragraph{Foundation of the Liberal Disorder}


To move down to the horizontal level of the issues, the whole left-right diremption arose in the aftermath of the French Revolution. The Revolution overthrew the temporal power and spiritual authority at that time, not just contingently presuming they were corrupt and needed replacement, but actually in principle. That has been the modus operandi of all succeeding revolutions, which have attacked the divine order in myriad ways. They cannot therefore be recognized without a deep understanding of the divine and natural orders. Those who claim to be on the right must belong to the streams of thought than began with \textbf{Edmund Burke} and \textbf{Joseph de Maistre}. Anything else is most likely arbitrary and unrelated to any higher order. Burke summarizes the foundations of the modern liberal order thusly:

\begin{quotex}
Instead of the religion and the law by which they were in a great politic communion with the Christian world, they have constructed their Republic on three bases, all fundamentally opposite to those on which the communities of Europe are built. Its foundation is laid in Regicide, in Jacobinism, and in Atheism; and it has joined to those principles, a body of systematic manners which secures their operation.
\flright{\textsc{Edmund Burke}, \emph{Letters on a Regicide Peace}}
\end{quotex}


\begin{itemize}


\item \textbf{Regicide}: all non-democratic governments are an usurpation.

\item \textbf{Jacobinism}: revolt against property. Distributing the wealth among the people.

\item \textbf{Atheism}: denial of God as moral Governor of the world.

\end{itemize}


By ``manners'', Burke means something like the ``culture wars''. As such, manners are more important than laws in defining a society. Manners, therefore, count for more than intellectual arguments. In his TV show broadcast in 1955, Bishop Fulton J Sheen\footnote{\url{https://www.youtube.com/watch?v=w5iUYQ\_yF\_M}} discussed this in regard to Communism in Russia (whose eventually downfall he foresaw). He pointed out that if the communist atheists wanted to oppose religion on the basis of ideology, they would lose. That is why the cultured despisers of religion, like Voltaire, resort to ridicule rather than intellectual debate. Hence, those opposed to the modern world must oppose it in manners also. So called ``conservatives'', neo-reactionaries, etc., need to disdain the dominant culture: its entertainment, its education, its science, its secularism, its licentiousness, its morals, and so on. Otherwise, it is a waste of time and effort.

Note that a system of monarchy does not necessarily follow from Burke's critique. Rather, any system that opposes the democratic ideal of \emph{vox populi, vox dei} may be supportable. The opposite of Jacobinism is anti-socialism, and opposition to the excesses of liberty, equality, fraternity. We can't get into details right now, but we nevertheless have provided a plausible framework to evaluate public policy and positive law.

\flrightit{Posted on 2016-05-21 by Cologero}

\begin{center}* * *\end{center}

\begin{footnotesize}
\begin{sffamily}

\texttt{francismercuri on 2016-06-03 at 22:52 said:}

Re: ``Me's'' post vis a vis ``workers wages''.

You mentioned that ``economists make a point that such an increase would lead to unemployment''; while I might tend to agree, ``increase'' is not the exclusive remedy suggested above (in fact, no remedy at all was suggested, only the proposition that ``workers must be paid a reasonable or living wage'').

Consequently, we must find a ``traditional'' economic solution to the proposal--other than ``raising minimum wage''. What other solution could there be?

It isn't my intention to ``sell'' any specific economics herein; but, only indicate that there exist alternatives to explore with a ``traditional'' basis. Fundamentally, any attempt at economic justice at this point must be preceded by monetary reform; neither band-aids, nor socialism.

One such interesting and viable ``school'' of monetary reform is rooted in Catholic social teaching, and has developed as the ``Just Third Way''--of which Michael Greaney has worked tirelessly to promote and educate upon. Of all the ``options'' on the table, the JTW seems the most reasonable, logical, and most realistic to transition toward:
\url{http://www.cesj.org/learn/just-third-way/}


From an individual ``interior'' point of view, ``externals'' such as wealth and economic prosperity should NOT impact one's pursuit of spiritual development, virtue, or eudaimonia--as only virtue is good and can only produce eudaimonia.

However, external prosperity has ``value'', can assist in attaining virtue (when used appropriately), and is preferable to demographic privation. Individual economic prosperity moreover cast over a vast demographic distribution, naturally leads to more widespread private property ownership, and is thus desirable, as it, 1) preserves the function of operative ``free will'', and 2) is a bulwark contra socialism and leftist incursion. Having a mass class of perpetually angry ``workers'' on the other hand, is nothing but an invitation and welcome mat for the far left--so, some sort of reform MUST happen soon, and the discontent MUST be addressed--lest, these ``workers'' CONTINUE to drift to the enchantment of the far left--which is what is happening in ``real world'' reality. This conflict is precisely why the ``millennilals'' have gone far-left, and endorse a ``Bernie Sanders''.

We need to make the West ``Right Again'', and the only way to do that at the social level, is offer SOMETHING viable to the young masses, contra the left, offering them present, and future financial security, based on sound principles that have concrete results, here and now. Monetary reform--lest, the masses continue drifting left.

\hfill

\end{sffamily}
\end{footnotesize}

